\documentclass[../../../../main.tex]{subfiles}
\begin{document}

\begin{flushright}
\footnotesize\textit{【自動文字起こし・要確認】}
\end{flushright}

\noindent
\textbf{[解]} (1) $A_n = {}_m\mathrm{C}_n$ とする。$m = 2^k$ の時、
\[
A_n = \frac{(2^k)!}{n! (2^k - n)!}
\]
である。この分子は
\[
\left[ \frac{2^k}{2} \right] + \left[ \frac{2^k}{2^2} \right] + \dots + \left[ \frac{2^k}{2^k} \right] \quad \text{回} \quad (\equiv B \in \mathbb{Z}_{\ge 0})
\]
2で割り切れる。一方、$1 \le n \le m-1$ の時 分母は、
\[
\sum_{l=1}^k \left\{ \left[ \frac{n}{2^l} \right] + \left[ \frac{2^k - n}{2^l} \right] \right\} \quad \text{回} \quad (\equiv C \in \mathbb{Z}_{\ge 0})
\]
2で割り切れる。ここで、ガウス記号の性質から、
\[
\left[ \frac{n}{2^l} \right] + \left[ \frac{2^k - n}{2^l} \right] \le \left[ \frac{2^k}{2^l} \right]
\]
であり、$l = k$ の時明らかに等号不成立だから、$l$ について足して、
\[
C < B
\]
だから、分子の方が多く 2で割り切れるので、$A_n$ は偶数である。 \hfill $\qed$

\bigskip

\noindent
(2) まず $m \in \text{even}$ の時、${}_m\mathrm{C}_1 = m \in \text{even}$ で不適だから、以下 $m \in \text{odd}$ で考える。この時、
$t \in \mathbb{N}$ として、$m = t \cdot 2^k - 1$ とおける ($t \in \text{odd}$)。$k \ge 3$ の時、$r = 2^k$ とすると $2 \le r < m$ を満たし、
\begin{align*}
{}_m\mathrm{C}_r &= \frac{m!}{r!(m-r)!} = \frac{m!}{(r-1)!(m-r+1)!} \cdot \frac{m-r+1}{r} \\
&= \frac{2^k(t-1)}{2^k} \cdot {}_m\mathrm{C}_{r-1} = (t-1) \cdot {}_m\mathrm{C}_{r-1} \equiv 0 \pmod 2 \quad (\because t \in \text{odd}, {}_m\mathrm{C}_{r-1} \in \mathbb{Z})
\end{align*}
だから不適なので、$t=1$ が必要。逆に $t=1$ の時
\[
{}_m\mathrm{C}_n = {}_{2^k-1}\mathrm{C}_n = {}_{2^k}\mathrm{C}_n - {}_{2^k-1}\mathrm{C}_{n-1} \quad (n \ge 1)
\]
及び ${}_m\mathrm{C}_0 = 1$, (1)から ${}_{2^k}\mathrm{C}_n \in \text{even}$ だから、帰納的に ${}_m\mathrm{C}_n \in \text{odd}$ となり、十分。

以上から、$m = 2^k - 1 \; (k \in \mathbb{N})$ である。 \hfill $\qed$

\newpage

\end{document}
